\documentclass[11pt,a4paper]{article}
\usepackage[T1]{fontenc}
\usepackage[utf8]{inputenc}
\usepackage{lmodern}
\DeclareUnicodeCharacter{1E93}{\d{z}}
\usepackage{amsmath,amssymb}
\usepackage[margin=25mm]{geometry}
\usepackage[hidelinks]{hyperref}
\hypersetup{pdftitle={Variational lattice-gap bounds: plaquettes, electric trials and holonomies},pdfauthor={navstokgap research working notes}}
\newcommand{\tightlist}{\setlength{\itemsep}{0pt}\setlength{\parskip}{0pt}}
\setlength{\emergencystretch}{3em}
\setcounter{secnumdepth}{0}
\begin{document}
\hypertarget{variational-lattice-gap-bounds-plaquettes-electric-trials-and-holonomies}{%
\section{Variational lattice-gap bounds: plaquettes, electric trials and
holonomies}\label{variational-lattice-gap-bounds-plaquettes-electric-trials-and-holonomies}}

For a specified smooth gauge-invariant link function, the Feynman--Bijl
identity bounds the physical Kogut--Susskind gap above by its electric
gradient norm over its variance. Plaquettes give the strong-coupling
order; an averaged spatial Polyakov loop gives a complementary exact
susceptibility bound. Replacing its expectations by constant-mode values
is a separate unproved small-volume step.

An abelian electric trial gives a structure-factor gaplessness
criterion. For a \textbf{semisimple} gauge Lie algebra, a nonzero
electric trial with constant colour coefficients is excluded by Gauss's
law. This exclusion neither proves a positive gap nor a universal string
energy for all dressed trials. Continuum smearing and positive-measure
moments have their own
\href{finiteness-half-flowed-susceptibility.md}{canonical treatment}.

\hypertarget{the-feynmanbijl-inequality}{%
\subsection{The Feynman--Bijl
inequality}\label{the-feynmanbijl-inequality}}

Let \(H\) be self-adjoint and bounded below with simple ground state
\(\psi_0\) and next eigenvalue \(E_1\), and let \(O\) be symmetric with
\(\psi_0\in D(O)\), \(O\psi_0\in D(H)\) and
\(\langle\psi_0,O\psi_0\rangle=0\). Then \(O\psi_0\perp\psi_0\), so by
the min-max principle
\[E_1-E_0\ \le\ \frac{\langle O\psi_0,(H-E_0)O\psi_0\rangle}{\|O\psi_0\|^2}
=\frac{\tfrac12\langle\psi_0,[O,[H,O]]\psi_0\rangle}{\langle\psi_0,O^2\psi_0\rangle},\]
using \([O,[H,O]]=2OHO-O^2H-HO^2\) and \((H-E_0)\psi_0=0\). If \(O\)
commutes with the Gauss generators, \(O\psi_0\) lies in the physical
sector and the same bound holds for \(\Delta^{\rm phys}\). (Feynman,
Phys. Rev.~94 (1954) 262, metadata; the argument is the two lines
above.)

\hypertarget{link-functions-and-strong-coupling-order}{%
\subsection{Link functions and strong-coupling
order}\label{link-functions-and-strong-coupling-order}}

Take
\(H=\frac{\hbar c}{a}[\frac{g^2}{2}\sum_\ell(-\Delta_\ell)+\frac2{g^2}\sum_p(N-\operatorname{Re}\operatorname{tr}U_p)]\)
as in \href{mass-gap-obligations-lattice.md}{the obligations map} §1 and
a real function \(O=f(U)\) with \(\langle f\rangle=0\). Then
\([V,f]=0\), and for the Laplacian \(\Delta_\ell=\sum_aX_a^2\) in the
left-invariant fields \(X_a\) of link \(\ell\),
\[[f,[-\Delta_\ell,f]]=2\sum_a(X_af)^2=:2|\nabla_\ell f|^2,\] because
\([\Delta,f]h=(\Delta f)h+2\sum_a(X_af)(X_ah)\) and the second
commutator with \(f\) removes the \((\Delta f)\) terms and leaves
\(-2\sum_a(X_af)^2h\). Hence

\textbf{Proposition 1.} For every real gauge-invariant \(f(U)\) with
\(\langle f\rangle_0=0\),
\[\Delta^{\rm phys}_{a,L}\ \le\ \frac{\hbar cg^2}{2a}\;
\frac{\sum_\ell\langle|\nabla_\ell f|^2\rangle_0}{\langle f^2\rangle_0}.\]

\textbf{Corollary 2 (plaquette).} With
\(f=\operatorname{Re}\operatorname{tr}U_p-\langle\operatorname{Re}\operatorname{tr}U_p\rangle_0\):
for a link \(\ell\) of \(p\),
\(X_a\operatorname{tr}U_p=\operatorname{tr}(T^aM)\) up to a phase, with
\(M\) a cyclic rearrangement of \(U_p\), so by the \(SU(N)\)
completeness relation
\(\sum_a\operatorname{tr}(T^aM)\operatorname{tr}(T^aM^\dagger)=\tfrac12[\operatorname{tr}MM^\dagger-\tfrac1N|\operatorname{tr}M|^2]\le\tfrac N2\),
and
\(|\nabla_\ell\operatorname{Re}\operatorname{tr}U_p|^2\le|\nabla_\ell\operatorname{tr}U_p|^2\le N/2\).
Four links give \(\sum_\ell|\nabla_\ell f|^2\le2N\) and
\[\Delta^{\rm phys}_{a,L}\ \le\ \frac{\hbar c\,g^2N}{a\,\operatorname{Var}_0(\operatorname{Re}\operatorname{tr}U_p)} .\]

\textbf{Corollary 3 (two-sided strong-coupling order).} At \(g=\infty\)
the ground state is the Haar measure and
\(\operatorname{Var}(\operatorname{Re}\operatorname{tr}U_p)=\tfrac12\langle|\operatorname{tr}U|^2\rangle+\tfrac12\operatorname{Re}\langle(\operatorname{tr}U)^2\rangle\)
equals \(1\) for \(SU(2)\) and \(\tfrac12\) for \(SU(N\ge3)\), by
Peter--Weyl orthogonality of the fundamental character and the vanishing
of \(\int(\operatorname{tr}U)^2dU\) for \(N\ge3\). The ground-state
expectation of a local function is continuous in \(g^{-4}\) at
\(g=\infty\) uniformly in the volume (Yarotsky's Theorem 1, conclusion
4, analyticity of \(\omega(A)\) in the perturbation parameter; the
finite-volume statement is analytic perturbation theory of a simple
isolated eigenvalue), so there is \(g_1(N)\) with
\(\operatorname{Var}_0\ge\tfrac14\) for \(g\ge g_1\). Combined with
\href{strong-coupling-uniform-gap.md}{the strong-coupling note},
\[\gamma\,\frac{g^2}{2}\,\frac{N^2-1}{2N}\,\frac{\hbar c}{a}\ \le\ \Delta^{\rm phys}_{a,L}\ \le\ 4N\,g^2\,\frac{\hbar c}{a}
\qquad(g\ge\max(g_0,g_1)),\] uniformly in the lattice size. The
strong-coupling gap is \(g^2\hbar c/a\) times a number bounded on both
sides.

\textbf{Why the plaquette says nothing about the continuum.} As
\(a\to0\) along the scaling curve the links approach the identity and
\(\operatorname{Var}(\operatorname{Re}\operatorname{tr}U_p)\to0\) (the
plaquette angle fluctuates on the scale \(g\), and its square on the
scale \(g^4\)), so the right side of Corollary 2 grows without bound.
The plaquette is an ultraviolet observable whose excitation costs an
energy of order \(\hbar c/a\); the finiteness half \(m<\infty\) of the
conjecture requires an operator whose ratio
\(\sum_\ell\langle|\nabla_\ell O|^2\rangle/\operatorname{Var}O\) stays
of order \(a/\ell\) for a physical length \(\ell\). The
smeared-susceptibility treatment states the additional continuum
hypotheses.

\hypertarget{abelian-electric-trials-and-the-structure-factor}{%
\subsection{Abelian electric trials and the structure
factor}\label{abelian-electric-trials-and-the-structure-factor}}

For compact \(U(1)\) write \(U_\ell=e^{i\theta_\ell}\),
\(E_\ell=-i\partial/\partial\theta_\ell\),
\(\theta_p=\sum_{\ell\in\partial p}\epsilon_{p\ell}\theta_\ell\), and
\[H=\frac{\hbar c}{a}\Big[\frac{g^2}{2}\sum_\ell E_\ell^2+\frac1{g^2}\sum_p(1-\cos\theta_p)\Big]\]
(other conventions rescale \(g^2\)). Gauss's law generators
\(G_x=\sum_{\ell\ni x}\pm E_\ell\) commute with every \(E_\ell\), so

\textbf{Lemma 4.} \(O_\varphi=\sum_\ell\varphi_\ell E_\ell\) with real
\(\varphi\) is gauge invariant, and \(\langle O_\varphi\rangle_0=0\).

\emph{Proof.} The \(E_\ell\) commute among themselves. Charge
conjugation \(\theta\mapsto-\theta\) is a unitary symmetry of \(H\)
under which the unique ground state is invariant and
\(E_\ell\mapsto-E_\ell\). \(\square\)

With \([E_\ell,h(\theta)]=-i\partial_\ell h\) and the lattice curl
\((d\varphi)_p=\sum_{\ell\in\partial p}\epsilon_{p\ell}\varphi_\ell\):
\([O_\varphi,\cos\theta_p]=i(d\varphi)_p\sin\theta_p\),
\([O_\varphi,\sin\theta_p]=-i(d\varphi)_p\cos\theta_p\), hence
\[[O_\varphi,[H,O_\varphi]]=\frac{\hbar c}{ag^2}\sum_p(d\varphi)_p^2\cos\theta_p,\]
and the variational identity above gives
\[\Delta^{\rm phys}\le\frac{\hbar c}{2ag^2}\,\frac{\sum_p(d\varphi)_p^2\,c_0}{\langle O_\varphi^2\rangle_0},
\qquad c_0=\langle\cos\theta_p\rangle_0,\] where \(c_0\) is independent
of \(p\) by translation invariance of the unique ground state. To bound
\(\langle O_\varphi^2\rangle\) below use the uncertainty relation with
the gauge-invariant \(\Theta=\sum_p\psi_p\sin\theta_p\):
\(\langle O^2\rangle\langle\Theta^2\rangle\ge\tfrac14|\langle[O,\Theta]\rangle|^2\)
with \(\langle[O_\varphi,\Theta]\rangle=-ic_0\sum_p\psi_p(d\varphi)_p\).
Choosing \(\psi=d\varphi\):

\textbf{Theorem 5.} If \(c_0>0\), then for every real \(\varphi\) on the
links with \(d\varphi\ne0\),
\[\Delta^{\rm phys}_{a,L}\ \le\ \frac{2\hbar c}{a\,g^2\,c_0}\;
\frac{\langle\Theta_\varphi^2\rangle_0}{\sum_p(d\varphi)_p^2},
\qquad\Theta_\varphi=\sum_p(d\varphi)_p\sin\theta_p .\] For a plane wave
\(\varphi_\ell=\cos(k\cdot x_\ell)\,u_{i(\ell)}\) the ratio is the
structure factor \(\hat S(k)\) of the plaquette sine at momentum \(k\),
and \(\Delta^{\rm phys}\le2\hbar c\,\hat S(k)/(ag^2c_0)\) for every
\(k\ne0\) allowed by the box.

\emph{Proof.} Insert the uncertainty bound into the previous display:
\(\langle O^2\rangle\ge c_0^2(\sum_p(d\varphi)_p^2)^2/(4\langle\Theta^2\rangle)\).
\(\square\)

\textbf{Reading.} In a Coulomb phase \(\sin\theta_p\) is a curl of a
nearly Gaussian field and \(\hat S(k)\to0\) as \(k\to0\); the theorem
then forces \(\Delta\to0\) with the box, which is the gaplessness Guth
and Fröhlich--Spencer prove at weak coupling (abstract, B78), obtained
here as a consequence of two ground-state quantities, \(c_0\) and
\(\hat S\). In a confining phase \(\hat S(0)>0\) and the bound is
finite, consistent with the strong-coupling gap. The theorem does not
decide which phase occurs; it shows that vanishing of this structure
factor is a sufficient criterion for gaplessness in the specified
channel. A nonvanishing upper bound gives no converse lower-gap
criterion.

\hypertarget{semisimple-gauss-law-and-the-limits-of-the-dressing-argument}{%
\subsection{Semisimple Gauss law and the limits of the dressing
argument}\label{semisimple-gauss-law-and-the-limits-of-the-dressing-argument}}

\textbf{Proposition 6.} Let \(G\) have semisimple Lie algebra and
\(O=\sum_\ell\sum_a\varphi^a_\ell E^a_\ell\) with c-number coefficients.
If \(O\) commutes with all Gauss generators, then \(\varphi=0\).

\emph{Proof.} Write \(E^a_\ell=L^a_\ell\) for the left generators of
link \(\ell\), which satisfy \([L^b_\ell,L^a_\ell]=if^{bac}L^c_\ell\),
and let \(R^b_\ell\) be the right generators, which commute with all
\(L^a_\ell\). The Gauss generator at \(x\) is
\(G^b_x=\sum_{s(\ell)=x}L^b_\ell+\sum_{t(\ell)=x}R^b_\ell\) up to the
sign convention. Then
\([G^b_x,O]=i\sum_{s(\ell)=x}f^{bac}\varphi^a_\ell L^c_\ell\), with no
contribution from the entering links. The \(L^c_\ell\) are linearly
independent operators, so \([G^b_x,O]=0\) for all \(b,x\) forces
\(f^{bac}\varphi^a_\ell=0\) for all \(b,c\) and every link, that is
\([\varphi_\ell,T^b]=0\) for all \(b\): \(\varphi_\ell\) lies in the
centre of the Lie algebra, which is zero under the semisimplicity
hypothesis. \(\square\)

The proof identifies the exact boundary: coefficients must lie in the
Lie algebra centre. A nonabelian group with a \(U(1)\) factor still
admits central abelian trials; the former unrestricted ``nonabelian''
statement was too broad.

Field-dependent gauge-covariant coefficients can produce invariant
electric trials, with ordering/domain obligations. A Wilson-loop
dressing is one example; local curvature coefficients are another. For
such an operator, \([K,O]\) generally contributes to the double
commutator, but \(O\) is a differential operator rather than a
multiplication function. Proposition 1's gradient formula cannot be
substituted for that extra term. The former universal positive cost
proportional to \(g^2|C|\) was not established. Gauge invariance removes
the specified c-number semisimple trial; it proves neither a mandatory
long string nor a nonvanishing dressing energy or lower gap.

\hypertarget{spatial-polyakov-loop-trials}{%
\subsection{Spatial Polyakov-loop
trials}\label{spatial-polyakov-loop-trials}}

\hypertarget{the-averaged-holonomy-and-its-exact-bound}{%
\subsubsection{The averaged holonomy and its exact
bound}\label{the-averaged-holonomy-and-its-exact-bound}}

Fix a direction \(i\). The lattice decomposes into \(N_s^2\) lines
parallel to \(i\), each of \(N_s\) links closing around the torus. For a
line \(\lambda\) let
\(P_\lambda=\operatorname{Re}\operatorname{tr}\prod_{\ell\in\lambda}U_\ell\),
the ordered product along the line; it is gauge invariant (a closed
loop), and \(\bar P_i=N_s^{-2}\sum_\lambda P_\lambda\). Let
\(O=\bar P_i-\langle\bar P_i\rangle_0\).

\textbf{Polyakov-loop bound.} For every \(g>0\), \(N_s\ge2\) and
\(SU(N)\) with the normalization of
\href{mass-gap-obligations-lattice.md}{the obligations map},
\[\Delta^{\rm phys}_{a,L}\ \le\ \frac{\hbar cg^2}{2a}\,
\frac{\big\langle\sum_\ell|\nabla_\ell\bar P_i|^2\big\rangle_0}{\operatorname{Var}_0\bar P_i}
\ \le\ \frac{\hbar c\,g^2N}{4L\,\operatorname{Var}_0\bar P_i}.\]

\emph{Proof.} \(O\) is a real gauge-invariant function of the links with
zero mean, so Proposition 1 of \href{lattice-gap-upper-bounds.md}{the
upper-bound note} applies and gives the first inequality. Each link in
direction \(i\) belongs to exactly one line and links in other
directions do not enter \(\bar P_i\), so
\(\nabla_\ell\bar P_i=N_s^{-2}\nabla_\ell P_{\lambda(\ell)}\). For a
link of \(\lambda\),
\(X_a\operatorname{tr}\prod U=\operatorname{tr}(T^aM)\) up to a phase
with \(M\) a cyclic rearrangement of the product, and the completeness
relation gives
\(\sum_a|\operatorname{tr}(T^aM)|^2=\tfrac12(N-|\operatorname{tr}M|^2/N)\le N/2\);
the real part has smaller gradient. Hence
\(\sum_\ell|\nabla_\ell\bar P_i|^2\le N_s^{-4}\cdot N_s^2\cdot N_s\cdot N/2=N/(2N_s)\),
and \(\frac{\hbar cg^2}{2a}\cdot\frac{N}{2N_s}=\frac{\hbar cg^2N}{4L}\).
\(\square\)

\hypertarget{conditional-small-volume-comparison}{%
\subsubsection{Conditional small-volume
comparison}\label{conditional-small-volume-comparison}}

For the formal weak-coupling constant-mode comparison, assume the links
are near the identity and the gauge-invariant content of a Polyakov line
is the holonomy of the constant mode:
\(\prod_{\ell\in\lambda}U_\ell\approx\exp(iLa_i)\) with \(a_i=a_i^bT^b\)
the constant mode of \href{low-dimensional-mass-gap.md}{G07} Proposition
7, whose ground-state width is \(|a_i|\sim g^{2/3}/L\) in natural units.
Since \(\operatorname{tr}T^b=0\),
\[P_\lambda\approx\operatorname{Re}\operatorname{tr}e^{iLa_i}=N-\tfrac14L^2|\vec a_i|^2+O(L^4|a|^4),\qquad
\nabla_\ell P_\lambda\approx\operatorname{Im}\operatorname{tr}(T^aLa_i)=\tfrac12La_i^a ,\]
so the variance of \(\bar P_i\) is second order and the gradient first
order in the constant mode:
\[\operatorname{Var}_0\bar P_i\approx\tfrac1{16}L^4\operatorname{Var}(|\vec a_i|^2),\qquad
\Big\langle\sum_\ell|\nabla_\ell\bar P_i|^2\Big\rangle_0\approx N_s^{-4}\cdot N_s^3\cdot\tfrac14L^2\langle|\vec a_i|^2\rangle
=\frac{L^2\langle|\vec a_i|^2\rangle}{4N_s}.\] For the \(SU(2)\)
constant-mode model of C133, \(La_i=g^{2/3}\xi_i\) with \(\xi_i\)
distributed by the unit ground state \(\psi_0^{(3)}\) of
\(\tfrac12(-\Delta_{\mathbb R^9}+\sum_{i<j}|\vec\xi_i\times\vec\xi_j|^2)\):
\[\frac{\hbar cg^2}{2a}\cdot\frac{g^{4/3}\langle|\vec\xi_i|^2\rangle/(4N_s)}{g^{8/3}\operatorname{Var}(|\vec\xi_i|^2)/16}
=\frac{\hbar c}{L}\,g^{2/3}\cdot\frac{2\langle|\vec\xi_i|^2\rangle}{\operatorname{Var}(|\vec\xi_i|^2)} .\]
The bound is therefore \((\hbar c/L)g^{2/3}\) times the pure number
\(2\langle|\vec\xi_i|^2\rangle/\operatorname{Var}(|\vec\xi_i|^2)\),
finite because the unit ground state decays faster than any power along
the valleys (the confining bound of C133 gives exponential decay of
\(\psi_0\) in \(\sum_i|\vec\xi_i|\)). Consistency with the exact
zero-mode gap \(\delta_1^{(3)}g^{2/3}\hbar c/L\) requires this number to
be at least \(\delta_1^{(3)}\), which is the Feynman--Bijl inequality
applied inside the zero-mode model with the trial operator
\(|\vec\xi_i|^2\); the model version of the Polyakov-loop bound is exact
and says exactly that.

What is proved and what is not: the Polyakov-loop bound is exact. The
two approximations, replacing the ground-state expectations by their
constant-mode values, are the leading order of the small-volume
expansion (Lüscher; abstract, B78) and are not proved here. The
statement that follows from them is that the exact lattice inequality
carries the \(g^{2/3}\) scaling of the small-volume corner on its upper
side; a lower bound with that scaling remains the open weak-coupling
target of STATE.

\hypertarget{infinite-volume-susceptibility-condition}{%
\subsubsection{Infinite-volume susceptibility
condition}\label{infinite-volume-susceptibility-condition}}

Suppose the lattice theory at fixed \(a\) and \(g\) has infinite-volume
gap \(\Delta_a=\liminf_{N_s\to\infty}\Delta_{a,L}>0\) (T2\('\)), and
write \(mc^2=\Delta_a\). the Polyakov-loop bound gives, for every \(L\),
\[\operatorname{Var}_0\bar P_i\ \le\ \frac{\hbar cg^2N}{4L\,\Delta_{a,L}} .\]
By translation invariance in the transverse plane,
\(\operatorname{Var}_0\bar P_i=N_s^{-2}\sum_{y}C_i(y)\) with
\(C_i(y)=\langle\operatorname{Re}P_\lambda(x)\operatorname{Re}P_\lambda(x+y)\rangle_c\)
the connected two-point function of the spatial Polyakov loop at
transverse separation \(y\). Hence

\textbf{Susceptibility consequence.} If the infinite-volume gap is
\(mc^2>0\), then along any sequence of boxes on which
\(\Delta_{a,L}\ge mc^2/2\),
\[\sum_{y\in\text{transverse plane}}a^2\,C_i(y)\ \le\ \frac{\hbar cg^2N}{2mc^2}\,L ,\]
and the averaged spatial Polyakov loop self-averages at least like
\(1/L\).

Reading: in a confining phase the spatial Polyakov loop of length \(L\)
has connected correlations decaying on a transverse scale \(\xi\), and
the left side is of order \(\xi^2\langle P^2\rangle_c\), so the
corollary is satisfied with room; in a phase where the loop orders or
its correlations do not decay, the left side grows like \(L^2\) and the
corollary forbids a gap. It is a necessary condition only, but it is one
written entirely in terms of a single gauge-invariant correlation and
the coupling, with no reference to the spectrum.

\hypertarget{the-haar-strong-coupling-comparison}{%
\subsubsection{The Haar strong-coupling
comparison}\label{the-haar-strong-coupling-comparison}}

At \(g=\infty\) the links are independent Haar variables, distinct lines
are independent, and
\(\langle|\operatorname{tr}U_1\cdots U_{N_s}|^2\rangle=1\) for a product
of independent Haar elements of \(SU(N)\) (the product is again Haar
distributed), so \(\operatorname{Var}P_\lambda=\tfrac12\) for \(N\ge3\)
and \(\operatorname{Var}\bar P_i=1/(2N_s^2)\). the Polyakov-loop bound
then reads \(\Delta\le\hbar cg^2NN_s^2/(2L)\), which grows with the
lattice and says nothing. The averaged loop is the wrong trial operator
in the disordered phase, as the plaquette is in the ordered one; the two
bounds are complementary, and neither is uniform in \(a\) along the
scaling curve.

\hypertarget{consequence-for-state}{%
\subsection{Consequence for STATE}\label{consequence-for-state}}

The exact upper bounds retain their specified trial states, groups,
variances, units and finite-volume domains. Strong-coupling order, a
vanishing abelian structure factor and a Polyakov susceptibility
constraint answer different questions. Neither a nonzero upper bound nor
the absence of one constant-colour trial establishes a lower gap. The
small-volume spectral transport and the continuum equal-time smearing
hypotheses remain open.
\end{document}
